\documentclass[10pt,a4paper]{moderncv}

%-------------------------------------------------------------------------------
% CONFIGURATION
%-------------------------------------------------------------------------------
\moderncvstyle[nosymbols]{classic}
\moderncvcolor{orange}
\nopagenumbers{}

\usepackage[scale=0.78]{geometry}
\setlength{\hintscolumnwidth}{3.3cm}
\usepackage{fontspec}
\usepackage{xeCJK}
\setmainfont{HaranoAjiGothic-Regular.otf}[
  BoldFont=HaranoAjiGothic-Bold.otf,
  ItalicFont=HaranoAjiGothic-Regular.otf,
  ItalicFeatures={FakeSlant=0.2},
  BoldItalicFont=HaranoAjiGothic-Bold.otf,
  BoldItalicFeatures={FakeSlant=0.2}
]
\setsansfont{HaranoAjiGothic-Regular.otf}[
  BoldFont=HaranoAjiGothic-Bold.otf,
  ItalicFont=HaranoAjiGothic-Regular.otf,
  ItalicFeatures={FakeSlant=0.2},
  BoldItalicFont=HaranoAjiGothic-Bold.otf,
  BoldItalicFeatures={FakeSlant=0.2}
]
\setCJKmainfont{HaranoAjiGothic-Regular.otf}[
  BoldFont=HaranoAjiGothic-Bold.otf,
  ItalicFont=HaranoAjiGothic-Regular.otf,
  ItalicFeatures={FakeSlant=0.2},
  BoldItalicFont=HaranoAjiGothic-Bold.otf,
  BoldItalicFeatures={FakeSlant=0.2}
]
\setCJKsansfont{HaranoAjiGothic-Regular.otf}[
  BoldFont=HaranoAjiGothic-Bold.otf,
  ItalicFont=HaranoAjiGothic-Regular.otf,
  ItalicFeatures={FakeSlant=0.2},
  BoldItalicFont=HaranoAjiGothic-Bold.otf,
  BoldItalicFeatures={FakeSlant=0.2}
]
\renewcommand{\familydefault}{\sfdefault}

%-------------------------------------------------------------------------------
% PERSONAL DATA
%-------------------------------------------------------------------------------
\renewcommand*{\addressfont}{\small\mdseries\raggedleft}

\name{新井}{理人}
\title{シニアソフトウェアエンジニア}
\address{}{長野県、日本}
\phone[mobile]{+81 (90) 7948-0077}
\email{licht.e.jima@gmail.com}
\social[linkedin]{rihito}
\social[github]{rhty}

\begin{document}

%-------------------------------------------------------------------------------
% NAME AND TITLE
%-------------------------------------------------------------------------------
\makecvtitle

%-------------------------------------------------------------------------------
% SUMMARY
%-------------------------------------------------------------------------------
\section{職務要約}
\textbf{GoとTypeScript}を中心に、バックエンド基盤およびAIプロダクトの開発に\textbf{7年以上}携わるシニアソフトウェアエンジニア／CTO。
LLMパイプライン、バックエンドサービス、クラウドインフラ、管理ツールを横断し、プロダクトの設計から提供までをリードしてきた。
マルチテナント型アクセス制御の設計、レガシーシステムのテスト・リリースプロセス改善、要件整理から本番リリースまでの一貫した推進を得意とする。
日本とカナダでの実務経験があり、日本語・英語を用いてビジネス部門と開発チームをつなぐことができる。

%-------------------------------------------------------------------------------
% AVAILABILITY
%-------------------------------------------------------------------------------
\section{希望条件}
2026年8月よりフルリモート案件に参画可能。
フルタイムはプロジェクト単位で最長6か月程度、パートタイムは週25時間程度で長期案件にも対応可能。
英語を使用する海外・国際チームの案件を特に希望。

%-------------------------------------------------------------------------------
% KEY SKILLS
%-------------------------------------------------------------------------------
\section{主要スキル}
\cvitem{バックエンド・言語}{Go、TypeScript（NestJS）、Ruby、Python、Google Apps Script}
\cvitem{AI・データ}{マルチモデルLLMパイプライン設計、OpenAI / Anthropic / Google API、データウェアハウス向けパイプライン、AIを活用したデータ処理}
\cvitem{フロントエンド}{React、Angular、shadcn/ui、MUI}
\cvitem{インフラ・DevOps}{AWS、GCP（Cloud Run、GKE、Cloud SQL、Pub/Sub、Cloud Build、Firebase）、Docker、Kubernetes、Terraform、GitHub Actions、GitLab CI/CD、Jenkins、Datadog、Sentry、Fastly、Cloudflare}
\cvitem{データベース}{PostgreSQL、MySQL、MSSQL、Redis、DynamoDB}
\cvitem{設計・セキュリティ}{マイクロサービス、マルチテナンシー、ACL / RBAC、機微な個人情報の暗号化・法令対応}
\cvitem{API・外部連携}{REST、OpenAPI / Swagger、gRPC、GraphQL、OAuth 2.0、Stripe、Ory Keto / Kratos}
\cvitem{品質・リード}{技術リード、CI/CD、レガシーコード改善、テスト戦略、コードレビュー、スクラムマスター}
\cvitem{言語}{日本語（母語）、英語（ビジネス：実務経験3年以上）、スペイン語（初級）}

%-------------------------------------------------------------------------------
% EXPERIENCE
%-------------------------------------------------------------------------------
\newpage
\section{職務経歴}

% -- VisionCompass
\cventry
    {2026年4月 -- 現在}
    {CTO}
    {株式会社VisionCompass}
    {日本（リモート）}
    {}
    {
      \begin{itemize}
        \item 2026年7月13日にリリースしたAIジャーナリング・睡眠アプリ\textbf{「ねるぞう」}において、バックエンド、インフラ、AI、管理画面の開発を統括。
        \item OpenAI、Anthropic、Googleの各モデルを組み合わせ、ノード構成や条件分岐を柔軟に定義できるグラフ型LLMパイプラインを設計・実装。ジャーナルの要約、会話分析、思考パターンの整理に活用。
        \item GCP Cloud Run上にDocker化したGoバックエンドとSupabaseを構築し、React・shadcn/uiによる管理画面を開発。
        \item SupabaseをジョブストアとしてCloud Runワーカーがポーリングする非同期処理を実装。LLM APIの待ち時間を活用できるよう、ワーカー内の処理を並列化。
      \end{itemize}
    }

% -- COTEN
\cventry
    {2025年8月 -- 2026年7月}
    {ソフトウェアエンジニア（契約社員）}
    {株式会社COTEN}
    {日本（リモート）}
    {}
    {
      \begin{itemize}
        \item TypeScriptとGoogle Apps Scriptを用いたAI連携のGoogleスプレッドシートツールを設計・実装し、人手では現実的でない規模の世界史データ整備を自動化。
        \item DWH as a Serviceにおけるマルチテナンシー、組織モデル、ACLの設計・実装を主導。
        \item TypeScript・NestJSによるデータパイプラインおよびバックエンドAPIを設計・実装。
        \item バングラデシュ拠点のエンジニアチームをマネジメントし、React管理画面をMUIからshadcn/uiへ移行する際の技術設計をチームと共同で検討。
      \end{itemize}
    }

% -- CIBC
\cventry
    {2025年1月 -- 2025年10月}
    {シニア・アプリケーション・デベロッパー}
    {CIBC（Canadian Imperial Bank of Commerce）}
    {カナダ・トロント（ハイブリッド）}
    {}
    {
      \begin{itemize}
        \item Go、TypeScript / Angular、MSSQL、Windows Serverを用いた、行員向けコーポレートセキュリティ事案管理・評価システムを開発・安定化。
        \item リファクタリングと自動テストの工数を開発計画に組み込むよう関係者を説得し、エラーの少ない安定したリリースに貢献。
        \item Jenkinsパイプラインを改善してリリースフローを標準化し、デプロイ時の負荷とリスクを低減。
        \item レガシー／新システム間のAPIスキーマおよびリバースプロキシ連携を改善し、セキュリティ上の問題やデータ不整合も発見。
        \item 部門責任者から最も成功したプロジェクトとして評価された開発に貢献。若手エンジニア向けにバックエンド設計・CI/CDの知識共有も実施。
      \end{itemize}
    }

% -- WealthPark
\cventry
    {2021年11月 -- 2025年2月}
    {シニアソフトウェアエンジニア／テックリード代行}
    {WealthPark株式会社}
    {東京（リモート）}
    {}
    {
      \begin{itemize}
        \item Go（go-kit）・Ory Ketoによるアクセス制御マイクロサービスをAWS EKS上に構築。Python / Airflow / AWS Lambdaによる外部不動産データ連携パイプラインも開発。
        \item \textbf{オルタナティブ投資プロジェクト（2022年4月 -- 2025年2月）}：
        \begin{itemize}
          \item 立ち上げ時から参画し、曖昧な要件を整理しながらGoマイクロサービス・gRPCによる投資家登録、ファンド管理機能を実装。
          \item 4名のエンジニアをリードし、単体テストとコードレビューを導入。全体のテストカバレッジを30\%から約70\%、新機能を約90\%まで向上。
          \item マイナンバーを扱う機能で、暗号化と法令要件への対応を実施。
          \item 英語を使用する環境で、海外開発チームと日本のステークホルダー間のコミュニケーションを橋渡し。
        \end{itemize}
      \end{itemize}
    }

% -- WED
\newpage
\cventry
    {2019年6月 -- 2021年10月}
    {サーバーサイド／リサーチエンジニア・スクラムマスター}
    {WED株式会社}
    {東京}
    {}
    {
      \begin{itemize}
        \item GCP上でGo、Ruby on Rails、Pythonによるバックエンドサービスを開発・運用。
        \item 約10万MAUの家計管理サービスにおける重大な並行処理問題を解決し、ユーザー影響が発生する前にデータ不整合を防止。
        \item スクラムマスターとしてアジャイルセレモニー、コードレビュー指針、Lintツールを導入し、ボトルネックの解消とチームの自律性向上を推進。
      \end{itemize}
    }

% -- FiNC
\cventry
    {2018年10月 -- 2019年6月}
    {サーバーサイドエンジニア（インターン）}
    {株式会社FiNC Technologies}
    {東京}
    {}
    {
      \begin{itemize}
        \item ヘルスケア・フィットネス向けモバイルアプリのREST APIをRuby on Railsで開発。
        \item OAuth 2.0によるLINE・Facebookのソーシャルログインを実装。
        \item マイクロサービス設計、データベース設計、ユーザー認証の基礎を習得。
      \end{itemize}
    }

% -- Freelance
\cventry
    {2020年2月 -- 現在}
    {フリーランス／副業}
    {}
    {}
    {}
    {
      \begin{itemize}
        \item \textbf{企業向けLLMパイプライン・管理ツール（TypeScript、Ruby on Rails、Google Apps Script）}：スタートアップの社内LLMワークフローと専用管理ツールを改善。レガシースクリプトの整理とUIコンポーネント追加により、パイプラインの信頼性と開発者体験を向上。
        \item \textbf{瞑想アプリのバックエンド（Ruby on Rails、GCP）}：ユーザー管理・音声コンテンツ機能を単独で開発し、Terraformを用いてGCPへデプロイ。
      \end{itemize}
    }

%-------------------------------------------------------------------------------
% EDUCATION
%-------------------------------------------------------------------------------
\section{学歴}

\cventry
    {2021年9月}
    {修士（工学：電気工学・情報システム）}
    {東京大学}
    {東京}
    {}
    {
      \begin{itemize}
        \item レシートを対象としたエンティティマッチング手法を研究（Python、機械学習）。
      \end{itemize}
    }

\cventry
    {2017年3月}
    {学士（工学：応用物理学）}
    {北海道大学}
    {北海道}
    {}
    {
      \begin{itemize}
        \item 複雑ネットワーク理論と計算手法を研究（Fortran）。
      \end{itemize}
    }

\end{document}
